\section{Discussion}
\label{sec:discussion}

The per-iteration breakdown in Table~\ref{tab:iterations} concentrates 80\% of the total improvement (12.98 of 16.15 nDCG points) in the single transition from Approach~2 to Approach~3. That is from a frozen-feature ensemble to an end-to-end fine-tuned ViT-L driven by soft-label supervision. The asymmetry along the iteration axis is itself the central observation, and it can be unpacked along three connected lines: the low ceiling of inference-time refinements over frozen backbones, the dominant effect of the encoder-and-loss swap, and the diminishing returns of post-fine-tuning refinements that motivate the future-work direction.

Approach~2 fuses three pre-trained models that were never optimized for the soft-label setting: JPoSE and MMEN are triplet-supervised, MI-MM is binary multi-instance MaxMargin, and none sees the relevance matrix $R$ during training. Equal-weight averaging of their normalized scores recovers $+1.78$~nDCG on top of JPoSE, in line with the canonical observation that retrieval ensembles reduce variance without correcting bias~\cite{wortsman2022modelsoups}. Because the three models share the same failure mode (they ignore graded relevance), averaging them produces a more stable estimate of the same flawed objective rather than one that is closer to nDCG. Equal weights were used in the absence of a relevance-aligned validation signal and weighted averaging tuned on a held-out split could plausibly add a fraction of a point but is unlikely to overcome the bias ceiling. Graph-diffusion re-ranking adds another $+0.51$~nDCG at no training cost, consistent with magnitudes reported for diffusion in image retrieval~\cite{iscen2017diffusion}. A small $k=20$ keeps the affinity graph sparse and avoids over-smoothing, moderate $\alpha=0.3$ blends rather than replaces the original signal, and the symmetric construction in Eq.~\eqref{eq:rerank}---propagating through both visual and textual graphs---prevents either modality from dominating the result.

The single largest jump, $+12.98$~nDCG between iterations~3 and~4, is driven by three coupled changes that we deliberately bundled into a single iteration boundary by replacing pre-extracted features with an end-to-end ViT-L/14 encoder pre-trained on egocentric video~\cite{zhao2023avion,zhao2023lavila}, replacing binary triplet/MaxMargin objectives by the SMS loss~\cite{wang2024symmetricmultisimilaritylossepickitchens100} that consumes the soft-label matrix $R$ directly, and jointly fine-tuning both encoders rather than freezing them. Disentangling the relative contribution of each factor would require ablations beyond our compute budget. What is clear is that the gain dominates everything else by an order of magnitude. We attribute most of the effect to the combination of the strong egocentric backbone and the relevance-aware loss. The SMS objective in Eq.~\eqref{eq:sms} weights both positive and negative pairs by their soft relevance, which aligns the gradient signal with the official nDCG metric (Eq.~\eqref{eq:ndcg}) and is the only mechanism in our pipeline that directly optimises it. The size of the gain is consistent with the observation in~\cite{wray2021semantic} that graded-relevance objectives can reshape retrieval orderings well beyond what binary contrastive losses achieve.

Once the encoder has been fine-tuned end-to-end with SMS, however, further gains diminish quickly. TTA (\texttt{--flip --clip-length 32}) yields $+0.68$~nDCG with no training cost. Flip-averaging adapts the canonical recipe of~\cite{krizhevsky2012alexnet} to video and exploits the near-symmetry of kitchen scenes under horizontal reflection, while the longer clip length addresses the long tail of segment durations. Extending fine-tuning from 10 to 16 epochs adds a further $+0.20$~nDCG, suggesting the model is approaching the plateau of single-checkpoint training. This corresponds precisely to the regime in which model-averaging strategies tend to be most cost-effective and which motivates the future-work direction in Section~\ref{sec:future}. The practical implication of Table~\ref{tab:iterations} is that, under a tight compute budget, encoder fine-tuning with relevance-aware supervision is the highest-yield investment by a wide margin: re-ranking and TTA are essentially free and should be applied by default, ensemble fusion of frozen baselines is also cheap but quickly hits a ceiling, and longer training is subject to rapidly diminishing returns. By contrast, the encoder swap returned $+12.98$~nDCG for the cost of a single A100 day. This ordering is consistent with what both the SMS-Loss authors~\cite{wang2024symmetricmultisimilaritylossepickitchens100} and the CR-CLIP authors~\cite{he2025crclip} report on the same benchmark: the largest contribution to their final scores comes from the encoder choice and the relevance-aware loss, with subsequent refinements closing a smaller residual gap.
